\section*{Discussion}\label{sec:discussion}

\commentauth{Angel}{From Scott's comment, maybe we remove the first section
altogther (or 1 paragraph at most?). We use that space to move some of the allozyme
discussion here. Andy notes: We leave this here now and wait for review.}

\subsection*{Genome assembly in highly polymorphic species}

Improvements in the affordability, accessibility, and accuracy in sequencing
technologies have increased the generation of high quality reference genomes across
the tree of life. Over the last decade, we have seen the research outcomes of large
genome assembly consortiums, e.g., the Vertebrate Genome Project
\citep{rhie_towards_2021} and the Earth BioGenome Project
\citep{blaxter_earth_2025}, tackling the production of high quality genomes across
diverse taxa. We have also seen the generation of complete genomes, sequencing
whole chromosomes at a telomere-to-telomere scale \citep{nurk_complete_2022,
yoo_complete_2025}, as well as the sequencing and assembly of some of the largest
and most repetitive genomes known in nature \citep{campos-dominguez_genome_2025}.
Best practice protocols have also become available to applying these methods to new
taxa in a reproducible fashion \citep{lawniczak_best-practice_2025,li_genome_2024}.
\vspace{0.2cm}

Despite these advances, there appears to be a gap in knowledge in the analysis of
highly polymorphic organisms, such as barnacles. High heterozygosity presents
problems during the assembly process, e.g., hindering the differentiation between
paralogs and haplotypes, which often results in the inflation of the genome size and
a high proportion of duplicated sequences \citep{guiglielmoni_phasing_2024,
asalone_regional_2020,roach_purge_2018,mochizuki_practical_2023}. The most common
approach is to manage these haplotypes \textit{a posteriori}, purging them after
initial assembly \citep{guan_identifying_2020,roach_purge_2018}; however, while this
approach can be successful in generating haploid assemblies, as we have shown here,
it is not without its drawbacks. First, purging haplotypes could lead to a reduction
in the gene completeness in the main assembly, as primary haplotypes might be
incorrectly labelled as haplotigs, leading to their removal. Secondly, the ability
to identify and remove haplotypes is dependent on the properties of the assembly.
Sequences that are too fragmented and/or too divergent from one another are likely
harder to purge due to difficulties in mapping. Moreover, when performing phased
diploid assemblies, the excessive presence of haplotypes will interfere with
phasing, leading to a main haploid assembly that is a mosaic between the two
haplotypes and/or the poor quality of the secondary assembly.
\vspace{0.2cm}

These difficulties drive the need for the development of methods and analysis
best-practices focused on the assembly of highly heterozygous genomes, and for a
further exploration of how elevated heterozygosity affects \textit{de novo} assembly
parameters.
In particular, the assembly of phased diploid genomes has seen several improvements
in recent years, e.g.,
the direct purging of haplotigs and the integration of Hi-C data to phase haplotypes
without the need of parental genotypes in \textsc{hifiasm}
\citep{cheng_haplotype-resolved_2022},
and the integration of Hi-C and ultra-long read data for phased telomere-to-telomere
haplotypes in \textsc{Verkko2} \citep{antipov_verkko2_2025}.
Yet, it is presently unknown if these new methods are all equally suitable for highly
heterozygous genomes and further work is likely needed before we can obtain
haplotype-resolved genome assemblies for hyperpolymorphic organisms.
Despite these limitations, our work here, assembling a chromosome-level assembly for
a heterozygous barnacle species --- even if not fully phased --- is a step in the
right direction of generating high-quality genomic resources for these
highly diverse and difficult genomes.

\subsection*{Genome evolution in large populations}

Evolutionary processes at large demographic scales are expected to leave distinct
signatures in the genome, several of which are observed in our Pacific acorn
barnacle assembly.
For example, large populations are expected to harbor increased levels of segregating
polymorphism, the product of both the reduction in the stochastic loss due to
genetic drift and the increase in the population-scale mutation rate.
In accordance, we observe a barnacle genome with extreme levels of genetic diversity,
both at the individual and population level (Figure~\ref{fig:sup_kmers};
Table~\ref{tab:sup_popgen_stats}), a substantial proportion of which is located in
coding sequences.
\vspace{0.2cm}

Similarly, large populations experience an increase in the efficacy of purifying
selection, which is expected to reduce the levels of diversity and increase 
conservation across functional sequences.
We see these patterns reflected in the \bgla{} genome.
Putatively functional sites, particularly coding exons and tRNAs, show both increased
levels of sequence conservation (Figure~\ref{fig:phastcons}) and a reduction in
population-level genetic diversity (Figure~\ref{fig:feature_pi};
Table~\ref{tab:sup_feature_pi}) compared to less constrained features of the genome.
Notably, the rank ordering of the enrichment of conserved elements, with tRNAs and CDSs
as the most enriched, parallels findings in insects \citep{siepel_evolutionarily_2005},
suggesting that the targets of purifying selection are broadly conserved across
arthropods despite the much larger and more repeat-rich genome of \bgla{}.
Increased efficacy of purifying selection is also reflected in the results of the
McDonald-Kreitman test (Figure~\ref{fig:mkTest}), where we observe a relative decrease
in the proportion of non-synonymous polymorphism in a central Oregon \bgla{} population,
suggesting that weakly deleterious variants are more efficiently purged by selection
\citep{li_reverse_2008}.
Positive selection is also expected to increase in efficacy in large populations,
resulting in the fixation of a larger fraction of beneficial non-synonymous mutations.
We observe this pattern as well: in this central Oregon population, the majority
(54 to 61\%, depending on the metric of $\alpha$ used) of non-synonymous substitutions
reached fixation due to positive selection, a proportion higher than that observed in
diverse \textit{Drosophila} populations \citep{haller_asymptoticmk_2017}.
\vspace{0.2cm}

In large enough populations, selection should be efficacious enough that its effects
extend to synonymous sites, where substitutions are normally considered neutral.
This could result in a bias in the usage of synonymous codons, as different codons are
favored over others due to GC biases, tRNA abundance, and other factors
\citep{galtier_codon_2018}.
These effects are observed in the \bgla{} assembly in the form of reduction in the
effective number of codons seen across protein-coding genes (Figure~\ref{fig:codon_enc}).
On average, \bgla{} uses $\approx$ 43 out of 61 available codons, a greater reduction
than that observed in \textit{D. melanogaster} (Figure~\ref{fig:sup_dmel_enc}).
This strong bias in codon usage further reflects an increase in the efficacy of natural
selection due to the large size of Pacific acorn barnacle populations.
Moreover, these biases put into question the neutrality of at least some synonymous
substitutions in this system, which could have important ramifications for the study of
molecular evolution; however, the extent of these effects remain presently unknown and
should be the subject of future study.
\vspace{0.2cm}

In contrast, the distribution of transposable and repeat elements in the
\bgla{} assembly was unexpected, as we find a highly repetitive genome
composed of over 60\% repeats (Figure~\ref{fig:genome_pi};
Table~\ref{tab:sup_earlgrey_summary}).
Previous work has shown that purifying selection plays a key role in the accumulation
and fixation of transposable elements \citep{merel_relaxed_2025,langmuller_genomic_2023,
bourque_ten_2018}, following the expectation that repetitive sequences are, on average,
slightly deleterious.
Given the increased efficacy of selection we observed in the rest of the genome,
our expectation was to find a relatively small proportion of repeats.
However, the high repeat content observed in the genome does not necessarily
reflect the low efficacy of selection in these populations.
Some could be low-frequency, slightly deleterious variants in the process of being purged
from the population.
Other repeats might represent beneficial genetic variation, as transposable elements are
known to play roles in adaptive evolution \citep{schrader_impact_2019}.
For example, repeat elements can act as regulators of gene expression
\citep{trizzino_transposable_2017,rech_population-scale_2022}, and as such may provide
the regulatory plasticity required for survival in rapidly changing
intertidal environments.
Regardless of their role, disentangling the population-wide repeat landscape and its
fitness effects is likely impossible with a single haploid reference assembly.
Thus, this observation drives the need for larger comparative genomic studies across
barnacle populations before we can fully understand the dynamics of genomic repeats
and transposable elements in this and other large populations.

\subsection*{Molecular evolution in large populations: reconciling \dnds{} and the
McDonald-Kreitman test}\label{sec:discussion_dnds_mkt}

One of the more striking patterns in our results is the apparent discordance between
our two measures of selection on protein-coding genes.
The average pairwise \dnds{} between \bgla{} and \textit{B. crenatus} is 0.206,
indicating that the majority of coding sites are under purifying selection.
Yet the McDonald-Kreitman test reveals a seemingly opposite pattern: $\approx61\%$
of non-synonymous substitutions fixed between these species were driven by
positive selection ($\alpha_\text{asym}$ = 0.610).
At first glance, these results appear contradictory: how can coding sequences be
under strong purifying selection while simultaneously showing pervasive adaptive
evolution?
\vspace{0.2cm}

The resolution lies in the fact that \dnds{} and the McDonald-Kreitman test measure
fundamentally different quantities. The \dnds{} ratio averages across all
non-synonymous sites in a gene, the vast majority of which are under functional
constraint, resulting in a ratio well below one. In contrast, the McDonald-Kreitman
test partitions changes into polymorphism and divergence, using synonymous changes
as a neutral baseline to ask whether there is an \textit{excess} of non-synonymous
fixations relative to the neutral expectation. This test is therefore sensitive to
the subset of non-synonymous mutations that escape purifying selection and reach
fixation, even when most non-synonymous sites in a gene are constrained.
Moreover, pairwise \dnds{} is calculated from a single sequence per species and
therefore incorporates segregating polymorphisms into its divergence estimate. In
species with many weakly deleterious non-synonymous variants segregating at
appreciable frequencies, this can inflate the non-synonymous divergence and distort
estimates of selection \citep{li_reverse_2008}. By explicitly separating
polymorphism from fixation, the McDonald-Kreitman test avoids this conflation.
\vspace{0.2cm}

\cite{li_reverse_2008} compared the performance of \dnds{} and the McDonald-Kreitman
test across human, \textit{Drosophila}, and yeast genomes and found that the two
approaches identify distinct but correlated sets of genes as targets of selection.
Notably, the McDonald-Kreitman test was more powerful than \dnds{} only in
\textit{Drosophila}, the species with the highest polymorphism, where it
identified nearly three times as many genes under positive selection. In the
lower-diversity human and yeast datasets, the pattern was reversed. This finding
suggests that the statistical power of the McDonald-Kreitman test is greatest in
high-polymorphism species, making it particularly well suited for systems like
barnacles where segregating variation is abundant.
\vspace{0.2cm}

The pattern we observe in \bgla{} is thus consistent with theoretical expectations
for large populations: strong purifying selection efficiently removes the majority
of deleterious non-synonymous mutations, keeping \dnds{} low, while at the same time
the increased efficacy of positive selection drives a disproportionate fraction of
the remaining non-synonymous substitutions to fixation, producing a high
$\alpha$. This coupling of strong purifying and positive selection is a hallmark
of evolution in large populations and has been observed in \textit{Drosophila}, where
$\alpha_\text{asym}$ is approximately 0.57 \citep{messer_frequent_2013}.
Our estimate of $\alpha_\text{asym}$ of 0.610 in \bgla{} is comparable to
or slightly higher than that in \textit{Drosophila}, consistent with the expectation
that the even larger effective population size in barnacles further increases the
efficacy of both modes of selection.
\vspace{0.2cm}

A locus-level example of this complementarity between tests emerges at the malic
enzyme paralogs: the MK test flags \textit{Me.2} through an excess of
non-synonymous fixations ($\alpha = 0.719$), the signature of recurrent positive
selection, while the classic HKA test flags \textit{Me.1} --- albeit at nominal
significance ($p = 0.032$, uncorrected) --- through a deficit of silent polymorphism
relative to divergence, the signature of a recent selective sweep at linked neutral
sites. That the two paralogs are detected by different tests underscores that the
tempo and mode of adaptive evolution can differ substantially even between close copies
of the same gene.
Because the two copies of NADP-dependent malic enzyme share deep sequence homology but
differ in genomic context, this paralog-level asymmetry cautions against treating
classical ``allozyme loci'' as monolithic units in genome-scale reanalyses.
\vspace{0.2cm}

The \textit{Gpi} result provides a second connection between our genome-wide analyses and
the allozyme literature.
\cite{hedgecock_temporal_1994} identified \textit{Gpi} as one of several \bgla{} loci
showing significant microgeographic heterogeneity in allele frequencies, but allozyme-era
methods could not distinguish whether that heterogeneity reflected protein-level
selection, drift, or chaotic recruitment among sweepstakes-limited larval cohorts
\citep{hedgecock_does_1994}.
Our sequence-level data now show that \textit{Gpi} in \bgla{} has accumulated a
substantial excess of non-synonymous fixations relative to the neutral expectation,
consistent with a history of directional positive selection on the protein itself.
This pattern contrasts with prior work at the \textit{Gpi} ortholog in
\textit{S. balanoides}, where the MK test has not yielded a comparable signal
\citep{nunez_footprints_2020,rand_ecological_2002} --- illustrating that ``same allozyme
locus'' comparisons across congeneric barnacles can obscure real divergence in
selective history.
\textit{Gpi} encodes a glycolytic enzyme at the branch point between glycolysis and the
pentose phosphate pathway; protein-coding variation at \textit{Gpi} in other marine
invertebrates has been linked to thermal tolerance and metabolic flux under osmotic and
thermal stress \citep{flight_allozyme_2012}, providing a plausible functional axis along
which directional selection could act in Pacific intertidal populations.
\vspace{0.2cm}

At \textit{Mpi}, the modest excess of non-synonymous polymorphism (NI $>$ 1) is
suggestive of --- but does not establish --- a regime of balancing selection.
\cite{hedgecock_temporal_1994} documented significant allele-frequency heterogeneity at
\textit{Mpi} across a 1-m intertidal transect at Bodega Harbor, a pattern analogous to
that subsequently documented at the \textit{Mpi} ortholog in \textit{S. balanoides}
\citep{nunez_footprints_2020,rand_ecological_2002}, where alternative alleles are
beneficial at different intertidal habitats.
However, intertidal microhabitat is approximately uniform in our central Oregon sampling,
so we cannot directly test the habitat-based balancing-selection hypothesis developed in
\textit{S. balanoides}.
Future sampling of \bgla{} across its geographic range and across intertidal
microhabitats will be needed to determine whether balancing selection at \textit{Mpi} is
shared across these systems or reflects distinct evolutionary histories in Atlantic and
Pacific barnacles.

\subsection*{Barnacles as models for evolution at biological extremes}

Our work shows that a Pacific acorn barnacle population from central Oregon exhibits
genetic diversity of over 5\% (mean genome-wide $\pi$ of 0.0506).
These levels of genetic diversity are higher than those observed for other model
systems known for their relatively large population sizes.
For example, the genome-wide $\pi$ we observe in this population is an order of
magnitude higher than that observed in \textit{Drosophila} (Figure~\ref{fig:sup_dmel_pi})
\citep{nunez_footprints_2025,li_comparative_2025,pool_population_2012,
langley_genomic_2012}.
Even the $\pi$ observed at 0-fold degenerate sites, the least diverse and among the
most constrained sites in the genome, is still an order of magnitude higher than the
diversity observed genome-wide in these \textit{Drosophila} populations.
While elevated, this level of diversity is not unique to barnacles.
Similar magnitudes of genome-wide $\pi$ ($\approx10^{-2}$) have been observed
in other taxa, including corals \citep{meziere_corals_2025}, parasitic nematodes
\citep{stevens_nematodes_2023}, bivalves \citep{liu_bivalves_2025}, and mosquitoes
\citep{kent_aedes_2025,ag1000g}.
Moreover, as with other marine invertebrates, barnacle life history traits are likely
direct drivers of the elevated polymorphism observed, as factors such as large census
sizes, small propagule size, long larval dispersal, and broad geographic distribution
have been positively correlated with high genetic diversity in these
organisms \citep{leffler_revisiting_2012,romiguier_comparative_2014}.
\vspace{0.2cm}

The magnitude of diversity observed in \bgla{} solidly places them among
hyperpolymorphic populations, those exhibiting diversity above 5\%
\citep{cutter_molecular_2013}. This adds barnacles to a growing list of emerging
hyperpolymorphic model systems, which are key to studying the limits of biological
variation and molecular evolution in natural populations. Moreover, our observations
of extreme genetic diversity also have important implications for the study of species
boundaries. For example, recent work by~\cite{hart_genomic_2025} suggests that levels
of above 1\% neutral diversity in Eukaryotes likely serve as biological limits for the
definition of species. Barnacles and other hyperpolymorphic populations put these
putative cutoffs into question. We observe levels of genome-wide diversity of over 5\%
in just a very small portion of the biological range of the Pacific acorn barnacle.
Nucleotide diversity is even higher ($\approx8\%$) at 4-fold degenerate, and putatively
neutral, sites (Figure~\ref{fig:feature_pi}; Table~\ref{tab:sup_feature_pi}).
Given the species' broader geographical and ecological distribution, genetic
diversity is likely to increase when looking across the whole range. This
observation reiterates the importance of future studies of genetic diversity
of this species over its whole geographic range, providing some insights on
the mechanisms that maintain population and species boundaries in the presence
of extreme genetic diversity.
\vspace{0.2cm}

Beyond species boundaries, the extreme diversity observed in barnacles also speaks to
a long-standing question in population genetics: the relationship between effective
and census population sizes. From our results on diversity, we
observe that this barnacle population has an effective size of approximately $10^{6}$.
However, these populations also show densities of $10^{4}$ individuals per square meter
\citep{menge_recruitment_2000}. Based on this estimate, the census size of the Pacific
acorn barnacle population in central Oregon should be, conservatively, in the billions of
individuals. While the general disparity between the effective and census size estimates
has been highlighted for decades, i.e., Lewontin's Paradox \citep{lewontin_genetic_1974},
the mechanisms underlying the differences have been long debated
\citep{gillespie_role_1999,gillespie_genetic_2000,smith_hitch-hiking_1974,
kimura_neutral_1983} and are still the subjects of study \citep{buffalo_quantifying_2021}.
Our work on the Pacific acorn barnacle provides a unique opportunity to empirically
study this disparity and the biological processes that maintain it.

